\documentclass[11pt]{article}
\usepackage[margin=0.85in]{geometry}
\usepackage{amsmath,booktabs,xurl,hyperref}
\usepackage[T1]{fontenc}
\setlength{\emergencystretch}{2em}
\title{Interrupted Controller Rerun and the 100-ms Latency Gap}
\author{Pan Dynamics Collision Avoidance Research}
\date{October 1, 2026}
\begin{document}
\maketitle

\section{Scope and interruption}
The current worktree was frozen before execution, including the existing
compiled model and geometry kernels. Its base commit is
\path{29d3e41}, with the uncommitted controller changes identified in the
source manifest. The principal differences from the preceding encounter-range
campaign are a 50-ms default RK4 step instead of 5 ms, consistent interpolated
midpoints, and optional compiled numeric kernels. The 30-m encounter range,
6-mm configured safety margin, 50-ms hold, and five-second soft improvement
budget remain. No algorithm or configuration was edited for this rerun.

Fourteen original threat cases were scheduled, seven at each of 8 and 15 m/s.
The user stopped the campaign during its fifth case and requested analysis of
five-second calls against an intended 0.1-s latency. The campaign process was
terminated on October 1 at approximately 12:12:57 America/Chicago. Four runs
had exported traces; one was interrupted and nine were not started. Only one
of the four exported runs completed its recovery criterion. This is not a
completed fourteen-case validation.

After the stop, work was limited to source inspection and analysis of saved
data. A MATLAB script only loads MAT snapshots and exports stored quantities;
it executes no controller, optimizer or plant simulation. The prepared
failure-replay and broader regression scripts were not executed.

\section{Saved results}
The plant uses independent tight ODE45 integration, with 31 audit samples per
hold. A separate Python rectangle implementation checks the saved samples.
No road constraints, observer errors, noise, delay or random draws are added.
The scenario target continues its original known trajectory in the plant
and audit; it is not removed at the controller's encounter-range threshold.

\begin{center}\small
\begin{tabular}{lrrrrl}
\toprule
8-m/s scenario & Holds & End (s) & Gap (mm) & Max call (s) & Outcome \\
\midrule
turningCrossing & 23 & 1.15 & 1955.794 & 5.017589 & Domain error \\
brakingLead & 104 & 5.20 & 3.835 & 5.045531 & No feasible continuation \\
crossing & 19 & 0.95 & 3237.293 & 5.028670 & Domain error \\
headOn & 261 & 13.05 & 7.072 & 5.015045 & Recovery confirmed \\
\bottomrule
\end{tabular}\end{center}
All 407 returned holds have zero reported hard residual and zero PCBF safety
slack. All recorded executed prefixes have strictly positive sampled distance;
this does not establish that the three interrupted-by-error scenarios would
remain collision-free. The braking-lead prefix violates the configured 6-mm
physical replay buffer despite its nominal admission. The two domain errors
occur after about 3.6 and 3.3 ms, respectively; they are not five-second
optimization timeouts. BrakingLead's next call fails after 5.011932 s.

Recovery requires all five errors below
$[0.1\,\mathrm m,1^\circ,0.1\,\mathrm{m/s},0.05\,\mathrm{m/s},0.01\,\mathrm{rad/s}]$
for five seconds beginning no earlier than eight seconds. The head-on run
meets this at 13.05 s; its subsequent-call median is 2.511 ms, but its maximum
is 5.015045 s. Fast nominal holds therefore coexist with slow avoidance holds.
An early failure must not be counted as a speedup over a completed prior run.

\section{A slow safe call is different from a failed feasibility search}
The longest saved call is brakingLead frame 90 at simulation time 4.45 s.
It takes 5.045531 s with 101 prediction holds, including 93 completion holds,
and invokes numerical solvers 30 times across eight recorded sequential steps
and additional correction work. Every recorded step skips the primary safety
solve because a nominally feasible incumbent already exists. All evaluated
step candidates have zero hard and safety residuals.

The first seven steps attempt zero CLF relaxation and receive local
infeasibility flags, then solve for positive relaxation and attempt predictive
cost refinement. Nonlinear admission/correction accompanies these proposals.
The CLF relaxation improves from 31.984157 to 23.492844, but the call reaches
the five-second budget before satisfying its CLF completion rule. This is
repeated improvement of an already safe candidate, not five seconds without
any safe input.

The brakingLead call at 5.20 s is different: it returns
\texttt{noFeasibleContinuation} with \texttt{timeLimit}. Its previous saved
plan has 91 holds, zero hard/safety residuals and minimum predicted speed
4.907092 m/s. The measured next state differs from its predicted successor;
in particular lateral velocity differs by $7.43984\,10^{-4}$ m/s and yaw
rate by $-1.45098\,10^{-3}$ rad/s. Hence exact shifted-witness reuse is
unavailable and the complete nonlinear continuation requires readmission or
restoration. The saved error alone does not establish which constraint
prevented restoration. No per-stage trace from that failed call was exported,
and the exact active obstruction is therefore unresolved without further
diagnosis. Failure of the local search is not a proof of physical infeasibility.

\section{Why the current organization misses 0.1 s}
The default \texttt{timeLimitSeconds=5} permits the complete improvement loop
to use fifty times the requested 0.1 s. It is a soft deadline: an in-flight
factorization or another operation can overrun it. A faster kernel allows
more search within this budget and does not itself shorten a budget-limited
call.

Each sequential step can solve a zero-CLF conic problem, a positive-CLF
conic problem and a predictive-cost conic problem, followed by nonlinear
rollouts, endpoint corrections and correction QPs. The outer loop permits
24 sequential steps; a correction helper permits five QPs, with backtracking
and full trajectory reevaluation. These bounds are not a 100-ms computational
design. Changing early controls can require repairing every later input and
the small terminal core. The braking-lead first plan still contains 193 holds
(9.65 s) despite the encounter-range change.

There is also sensitivity of the feasible prediction. Before the domain
errors, turningCrossing's saved minimum predicted speed is 1.001618 m/s and
crossing's is 1.038516 m/s, close to the 1-m/s model floor. Their measured
speeds at failure are 4.705472 and 6.171439 m/s, respectively, so the measured
initial state itself is not below the floor. The stored prediction and actual
successor differ even with exact state observation, because the one-step RK4
prediction is not the ODE45 plant map. These observations identify fragile
nominal continuation and numerical-model mismatch as concerns, but do not
alone prove that the larger RK4 step caused every failure. No paired
integration-step campaign was run.

The design objective should remain safety followed by CLF improvement.
Returning at zero PCBF slack while ignoring a positive CLF slack would violate
that objective. Instead, a 100-ms design needs substantially less repeated
convex work and trajectory repair, a reliably transferable numerical
prediction, and a bounded improvement policy whose safety and dissipation
properties are established together. Directly minimizing CLF slack can avoid
the preliminary zero-slack feasibility attempt in the same local problem;
tail parameterization and correction coupling also deserve redesign. These
are proposed directions, not implemented changes or measured latency promises.
Simply replacing the five-second timeout with 0.1 s would cut search short
without establishing an admissible result by that deadline.

\section{Validation and reproducibility}
Eight compiled-kernel consistency tests passed before the campaign; their
random test streams use seeds 7 and 11. This verifies numeric consistency at
the tested points, not integration accuracy relative to the physical replay.
The larger controller suite was prepared but not run after the user stopped
experiments. The four saved traces were independently audited in Python;
the source snapshot and compiled binary hashes were checked unchanged.

Raw files, commands and the interruption record are under
\path{/home/zai/.cache/collisionAvoidance/controller-rerun-20261001}.
The repository bundle contains the compact partial summary, per-case and
sampled results, the slow safe-call trace, saved-state inspection and manifest.
No simulation was resumed after the stop, and no result is claimed for the
interrupted acceleratingHeadOn run or the nine unstarted cases.
\end{document}
